\section{Discussion and limitations}
\label{sec:discussion}

In homogeneous Case~H, rotating and elongating the characteristic-length tensor shifts the branches, changes directional gap widths, and redirects group velocity, while the complete-gap measure remains negative. Across the sampled steering grid, orientation moves the principal direction but leaves the peak deviation nearly unchanged. Orientation therefore controls directional dispersion and energy-flow direction in this case.

Case~C shows the effect of periodic impedance contrast, which produces a positive complete gap. The reported width falls from $2.5732$ on a $4\times4$ mesh to $1.9179$ on a $64\times64$ mesh and has not converged. The coarse-mesh value is resolution-specific; further refinement is needed before assigning a converged width. Case~H and Case~C use different material and discretization settings, so their gap values do not form a controlled parameter sweep.

The bilayer comparisons provide different levels of evidence. We check the classical AlN/BaTiO$_3$ calculation against the Rytov relation and compare it visually with the published plot. The two gradient-elastic cases are reproductions only; their source parameter conventions are unresolved, and numerical comparison data are unavailable (Supplementary Material~S2).

The analysis is linear, local, isothermal, and two-dimensional under plane-strain assumptions. It does not include nonlinear gradient effects, dissipation, piezoelectric coupling, or three-dimensional polarization modes. The finite high-wavenumber phase-speed result applies only to the transverse branch and constitutive model examined here.
